% Standalone problem document, generated from the adjacent README.md.
% Compile with XeLaTeX. Mathematical content is maintained in Markdown.
\documentclass[11pt,a4paper]{article}
\usepackage[fontsize=10.5pt]{scrextend}
\usepackage[margin=22mm,headheight=15pt,headsep=9mm,footskip=11mm]{geometry}
\usepackage{fontspec}
\setmainfont{texgyrepagella}[Extension=.otf,UprightFont=*-regular,BoldFont=*-bold,ItalicFont=*-italic,BoldItalicFont=*-bolditalic]
\setsansfont{texgyreheros}[Extension=.otf,UprightFont=*-regular,BoldFont=*-bold,ItalicFont=*-italic,BoldItalicFont=*-bolditalic]
\setmonofont{texgyrecursor}[Extension=.otf,UprightFont=*-regular,BoldFont=*-bold,ItalicFont=*-italic,BoldItalicFont=*-bolditalic,Scale=MatchLowercase]
\usepackage{amsmath,amssymb}
\usepackage{unicode-math}
\setmathfont{texgyrepagella-math.otf}
\usepackage{xcolor,microtype,enumitem,needspace,fancyhdr}
\usepackage{bookmark}
\definecolor{ink}{HTML}{17344B}
\definecolor{accent}{HTML}{197D80}
\definecolor{muted}{HTML}{596773}
\hypersetup{colorlinks=true,linkcolor=accent,urlcolor=accent,pdftitle={RA-19 - Critical-point
count for corank-one approximation with a fixed
zero},pdfauthor={Open Problems in Numerical Linear Algebra}}
\urlstyle{same}
\setlength{\parindent}{0pt}
\setlength{\parskip}{5pt plus 1pt minus 1pt}
\linespread{1.04}
\setlength{\emergencystretch}{3em}
\widowpenalty=10000
\clubpenalty=10000
\displaywidowpenalty=10000
\allowdisplaybreaks[1]
\setlist{nosep,leftmargin=1.5em}
\providecommand{\tightlist}{\setlength{\itemsep}{0pt}\setlength{\parskip}{0pt}}
\providecommand{\pandocbounded}[1]{#1}
\pagestyle{fancy}
\fancyhf{}
\fancyhead[L]{\sffamily\footnotesize\color{muted}OPEN PROBLEMS IN NUMERICAL LINEAR ALGEBRA}
\fancyhead[R]{\sffamily\bfseries\footnotesize\color{accent}RA-19}
\fancyfoot[L]{\parbox[b]{0.9\textwidth}{\sffamily\fontsize{8}{10}\selectfont\color{muted}Editorial ratings; a bounded search cannot certify that a problem remains unsolved.\\Literature check: 2026-09-11}}
\fancyfoot[R]{\sffamily\footnotesize\color{muted}\thepage}
\renewcommand{\headrulewidth}{0.3pt}
\renewcommand{\headrule}{\hbox to\headwidth{\color{accent}\leaders\hrule height \headrulewidth\hfill}}
\makeatletter
\renewcommand{\section}{\Needspace{5\baselineskip}\@startsection{section}{1}{0pt}{12pt plus 2pt minus 2pt}{4pt}{\sffamily\bfseries\large\color{ink}}}
\renewcommand{\subsection}{\Needspace{4\baselineskip}\@startsection{subsection}{2}{0pt}{9pt plus 1pt minus 1pt}{3pt}{\sffamily\bfseries\normalsize\color{ink}}}
\makeatother
\setcounter{secnumdepth}{0}
\begin{document}
{\sffamily\small\color{muted}Randomized and low-rank approximation\par}
\vspace{3pt}
{\raggedright\hyphenpenalty=10000\exhyphenpenalty=10000\sffamily\bfseries\fontsize{21}{24}\selectfont\color{ink}Critical-point
count for corank-one approximation with a fixed zero\par}
\vspace{7pt}
{\sffamily\small\textbf{Difficulty:} challenging\quad\textbf{Importance:} interesting
to specialist\par}
{\sffamily\small\bfseries\color{accent}Status: Open\par}
\vspace{4pt}
{\color{accent}\hrule height 0.8pt}
\vspace{6pt}
\textbf{Topic:} Structured low-rank approximation and Euclidean distance
degree

\textbf{Rating rationale:} The target is an exact formula for a
constrained approximation problem in every dimension. It would quantify
the algebraic complexity of enumerating candidate nearest singular
matrices with a prescribed zero.

\subsection{Statement}\label{statement}

For \(n\ge3\), let

\[V_n=\{X\in\mathbb C^{n\times n}:\det X=0,\ x_{11}=0\}.\]

For generic \(U\in\mathbb C^{n\times n}\), consider the bilinear squared
distance

\[d_U(X)=\sum_{i,j=1}^n(x_{ij}-u_{ij})^2.\]

Its critical points on the smooth locus of \(V_n\) satisfy
\(\mathrm d(d_U)_X[Z]=0\) for every tangent vector \(Z\in T_XV_n\).
Their finite number for generic \(U\), meaning outside a proper
algebraic exceptional set, is the \emph{Euclidean distance degree}
\(\operatorname{EDdeg}(V_n)\). There is no complex conjugation in this
definition; over the reals the objective is squared Frobenius distance.

\textbf{Conjecture (Kubjas--Sodomaco--Tsigaridas, Conjecture 5.1,
nontrivial dimensions).}

\[\operatorname{EDdeg}(V_n)=5n-7\qquad\text{for every }n\ge3.\]

The dimension restriction is explicit here: at \(n=2\), the equations
reduce to \(x_{11}=0\) and \(x_{12}x_{21}=0\), a union of two linear
spaces with ED degree two. The printed conjecture omits a lower bound on
\(n\), while its supporting table begins at three. That exceptional
endpoint is not asserted to satisfy the formula.

\subsection{Evidence and numerical
significance}\label{evidence-and-numerical-significance}

Table 2 of the source reports \(8,13,18,23,28,33,38,43\) for
\(n=3,\ldots,10\), using numerical critical-point computations checked
against symbolic ideal degrees. These finite computations support the
conjecture but do not establish the formula for arbitrary \(n\).

Truncated SVD solves unconstrained Frobenius approximation, but a fixed
zero changes its critical equations. This target concerns corank-one
approximation in general square matrices; the
\href{https://github.com/ajt60gaibb/OpenProblemsInNLA/blob/main/randomized-and-low-rank-approximation/RA-20/README.md}{symmetric
rank-two zero-pattern problem} has different geometry and rank
constraints.

\subsection{References and status
check}\label{references-and-status-check}

\begin{itemize}
\tightlist
\item
  K. Kubjas, L. Sodomaco and E. Tsigaridas, \emph{Exact solutions in
  low-rank approximation with zeros}, Linear Algebra and its
  Applications 641 (2022), 67--97.
  \href{https://doi.org/10.1016/j.laa.2022.01.021}{DOI};
  \href{https://arxiv.org/abs/2010.15636v2}{current author manuscript},
  29 January 2022. Conjecture 5.1 and Table 2, manuscript p.19; §2
  defines the critical-point and ED-degree conventions.
\end{itemize}

On 2026-09-11, checked the current arXiv version, the authors'
publication lists, and targeted title, conjecture-number, ED-degree,
proof and 2026 searches. No later proof or counterexample for \(n\ge3\)
was located. The formulation and table were checked in arXiv v2; a
separate publisher PDF was not obtained. The scope correction at \(n=2\)
is disclosed above. This is a bounded literature check.
\end{document}
